\section{9. Conclusion}\label{conclusion}

For two model-derived HgCdTe absorption curves reconstructed from a
published figure, the fitted optical intercept is not invariant to the
downstream reporting convention. The largest immediately controlled
sensitivity arises from the declared fit-domain and weighting rule: the
free-exponent intercept moves by as much as 5.71 and 5.44 meV across the
predeclared grid. Reconstruction choice with nominal membership fixed
produces smaller maxima of 0.69 and 0.81 meV, while the visible
two-point plateau changes a non-boundary fitted intercept by
approximately 0.12 meV when omitted.

The complete boundary-excluded fitted-model spans are 5.09 and 6.68 meV,
but energy-resolved residuals show that the candidates are not equally
compatible with the reconstructed traces. A line-envelope screen
identifies grossly incompatible forms, yet its retained span changes
under stricter coverage thresholds and is therefore not treated as a
primary interval. Fixed-absorption crossings answer a separate
operational question and are excluded from the fitted-intercept span.

The supported conclusion is methodological and conditional: for these
two reconstructed, model-derived spectra, the fitted HgCdTe optical
intercept is materially conditioned by extraction model and especially
by fit-domain endpoints and weighting convention. The available source
record does not identify a unique physical gap, a total experimental
uncertainty budget, or a sub-meV ranking of empirical gap equations.
Native spectra, optical-inversion covariance, specimen-level composition
uncertainty, and independent replication are required to move beyond
this case-study boundary.
